\documentclass[10pt,a4paper]{amsart}
\usepackage{amsmath,amssymb,amsthm,graphicx}
% For commutative diagrams we use Xy-pic
\usepackage[all]{xy}

% We use 2cell for 2-commutative diagrams.
\xyoption{2cell}
% Theorem environments.
%
\theoremstyle{plain}
\newtheorem{theorem}[subsection]{Theorem}
\newtheorem{proposition}[subsection]{Proposition}
\newtheorem{lemma}[subsection]{Lemma}

\theoremstyle{definition}
\newtheorem{definition}[subsection]{Definition}
\newtheorem{example}[subsection]{Example}
\newtheorem{exercise}[subsection]{Exercise}
\newtheorem{situation}[subsection]{Situation}

\theoremstyle{remark}
\newtheorem{remark}[subsection]{Remark}
\newtheorem{remarks}[subsection]{Remarks}

\numberwithin{equation}{subsection}

% Macros
%
\def\lim{\mathop{\mathrm{lim}}\nolimits}
\def\colim{\mathop{\mathrm{colim}}\nolimits}
\def\Spec{\mathop{\mathrm{Spec}}}
\def\Hom{\mathop{\mathrm{Hom}}\nolimits}
\def\Ext{\mathop{\mathrm{Ext}}\nolimits}
\def\SheafHom{\mathop{\mathcal{H}\!\mathit{om}}\nolimits}
\def\SheafExt{\mathop{\mathcal{E}\!\mathit{xt}}\nolimits}
\def\Sch{\mathit{Sch}}
\def\Mor{\mathop{\mathrm{Mor}}\nolimits}
\def\Ob{\mathop{\mathrm{Ob}}\nolimits}
\def\Sh{\mathop{\mathit{Sh}}\nolimits}
\def\NL{\mathop{N\!L}\nolimits}
\def\CH{\mathop{\mathrm{CH}}\nolimits}
\def\proetale{{pro\text{-}\acute{e}tale}}
\def\etale{{\acute{e}tale}}
\def\QCoh{\mathit{QCoh}}
\def\Ker{\mathop{\mathrm{Ker}}}
\def\Im{\mathop{\mathrm{Im}}}
\def\Coker{\mathop{\mathrm{Coker}}}
\def\Coim{\mathop{\mathrm{Coim}}}

% Boxtimes
%
\DeclareMathSymbol{\boxtimes}{\mathbin}{AMSa}{"02}

%
% Macros for moduli stacks/spaces
%
\def\QCohstack{\mathcal{QC}\!\mathit{oh}}
\def\Cohstack{\mathcal{C}\!\mathit{oh}}
\def\Spacesstack{\mathcal{S}\!\mathit{paces}}
\def\Quotfunctor{\mathrm{Quot}}
\def\Hilbfunctor{\mathrm{Hilb}}
\def\Curvesstack{\mathcal{C}\!\mathit{urves}}
\def\Polarizedstack{\mathcal{P}\!\mathit{olarized}}
\def\Complexesstack{\mathcal{C}\!\mathit{omplexes}}
% \Pic is the operator that assigns to X its picard group, usage \Pic(X)
% \Picardstack_{X/B} denotes the Picard stack of X over B
% \Picardfunctor_{X/B} denotes the Picard functor of X over B
\def\Pic{\mathop{\mathrm{Pic}}\nolimits}
\def\Picardstack{\mathcal{P}\!\mathit{ic}}
\def\Picardfunctor{\mathrm{Pic}}
\def\Deformationcategory{\mathcal{D}\!\mathit{ef}}
\title{A property of quasi-compact quasi-separated schemes stable under thickenings, affine limits and splittings descends to Spec of a field}
\author{The Stacks Project Authors}
\date{stacks-project:etale-cohomology.tex:local-snapshot-20261011; GNU FDL 1.2 or later}
\begin{document}
\maketitle
\noindent\textit{Editorial note.} All delivered mathematical text is the author's own, copied from the pinned Stacks Project chapter file etale-cohomology.tex; no statement or proof step was rewritten, reordered or removed.\par\medskip
\noindent\textit{Editorial note.} The counted proof is the author's complete argument: The counted proof is the author's complete argument for the reduction, stating explicitly how each of the three assumptions is used in turn.\par\medskip
\noindent\textit{Editorial note.} The label lemma-formal-argument defines the statement; the counted proof is the proof environment immediately following the lemma in the source.\par\medskip
\noindent\textit{Editorial note.} The preamble is the chapter's own native preamble excerpt already used for this source: the Xy-pic setup (native preamble lines 14-18) and the native theorem declarations with the mathematical macros (native preamble lines 154-217). The class block and the externaldocument lines of the native preamble are not included, because the assembler supplies the wrapper class and those external documents are not part of this material.\par\medskip
\noindent\textit{Allowed standard prerequisites (editorial).}\begin{itemize}
\item Thickenings and limits of schemes over a base T. Same chapter, the standard terminology of thickenings and of filtered limits with affine transition morphisms used in the three assumptions. Supplies the objects over which the property P is required to propagate.
\item Generic point of an irreducible component. Same chapter, the generic point eta of an irreducible component of a reduced quasi-compact quasi-separated scheme. Supplies the point eta = Spec(K) at which the conclusion is obtained, as a limit of affine open neighbourhoods V.
\end{itemize}\medskip
\begin{lemma}
\label{lemma-formal-argument}
Let $T$ be a quasi-compact and quasi-separated scheme.
Let $P$ be a property for quasi-compact and quasi-separated
schemes over $T$. Assume
\begin{enumerate}
\item If $T'' \to T'$ is a thickening of quasi-compact and
quasi-separated schemes over $T$, then $P(T'')$ if and only if $P(T')$.
\item If $T' = \lim T_i$ is a limit of an inverse system of
quasi-compact and quasi-separated schemes over $T$ with affine
transition morphisms and $P(T_i)$ holds for all $i$, then
$P(T')$ holds.
\item If $Z \subset T'$ is a closed subscheme with
quasi-compact complement $V \subset T'$ and $P(T')$ holds,
then either $P(V)$ or $P(Z)$ holds.
\end{enumerate}
Then $P(T)$ implies $P(\Spec(K))$ for some morphism $\Spec(K) \to T$
where $K$ is a field.
\end{lemma}

\begin{proof}
Consider the set $\mathfrak T$ of closed subschemes $T' \subset T$
such that $P(T')$. By assumption (2) this set has a minimal element,
say $T'$. By assumption (1) we see that $T'$ is reduced.
Let $\eta \in T'$ be the generic point of an irreducible
component of $T'$. Then $\eta = \Spec(K)$ for some field
$K$ and $\eta = \lim V$ where the limit is over the affine
open subschemes $V \subset T'$ containing $\eta$.
By assumption (3) and the minimality of $T'$ we see
that $P(V)$ holds for all these $V$. Hence $P(\eta)$
by (2) and the proof is complete.
\end{proof}

\end{document}
